\begin{abstract}
    \noindent The lattice diameter of a bounded set $S \subset \R^d$ measures the maximal number of lattice points in a segment whose endpoints are lattice points in $S$. Such a segment is called a lattice diameter segment of $S$. This simple invariant yields interesting applications and challenges. We describe a polynomial-time algorithm that computes lattice diameter segments of lattice polygons and show that computing lattice diameters of semi-algebraic sets in dimensions three and higher is $\NP$-hard.
    We prove that the function that counts lattice diameter segments in dilations of a lattice polygon is eventually a quasi-polynomial in the dilation factor.
    % in dimension two we show that the count of lattice diameter segments under dilation agrees with a quasi-polynomial.
    We also study the number of \mbox{directions that} lattice diameter segments can have. Finally, we prove a Borsuk-type theorem on the number of parts needed to partition a set of lattice points such that each part has strictly smaller \mbox{lattice diameter}.
\end{abstract}
